\documentclass[11pt,twoside]{article}
\usepackage[utf8]{inputenc}
\usepackage[margin=1in]{geometry}
\usepackage{titlesec}
\usepackage{parskip}
\usepackage{fancyhdr}
\usepackage{amsmath, amssymb}
\usepackage[round,sort&compress]{natbib}
\usepackage{hyperref}

% Approximation of JMLR style parameters
\setlength{\parindent}{0pt}
\setlength{\parskip}{0.3em}
\titlespacing*{\section}{0pt}{0.6em}{0.3em}
\titlespacing*{\subsection}{0pt}{0.5em}{0.2em}

\pagestyle{fancy}
\fancyhf{}
\lhead{\footnotesize STAT8101}
\chead{\footnotesize \href{http://www.cs.columbia.edu/~blei/teaching.html}{Invariance, Causality, and Applications}}
\rhead{\footnotesize \href{http://www.cs.columbia.edu/~blei/}{David M. Blei}}
\cfoot{\thepage}

\begin{document}
\vspace*{-1.5cm}
\begin{center}
{\Large \textbf{Reader Report: Week 12}} \\
\vspace{0.2em}
{\large Thompson Sampling, Latent-Space BO, and Causal Inference from Text} \\
\vspace{0.3em}
Daniel Hardesty Lewis (\texttt{dl3645}) $\cdot$ Spring 2026
\end{center}
\vspace{-1.5em}
\thispagestyle{fancy}

\section*{Summary}

\textit{BO completion.} Thompson sampling draws a posterior sample $\tilde{f} \sim p(f \mid \mathcal{D}_t)$ and maximizes it, achieving regret $O^*(\sqrt{T\gamma_T})$ comparable to GP-UCB while requiring no $\beta_t$ schedule and admitting natural batch parallelism via independent draws \citep{russo2014thompson,garnett2023bayesian}. Latent-space BO places a learned low-dimensional embedding between the raw search space and the GP, making Thompson sampling tractable where UCB acquisition optimization would otherwise degrade --- a learned analog of REMBO.

\textit{Causal inference from text.} \citet{egami2022texts} establish the core taxonomy: text appears as a \textit{treatment}, \textit{outcome}, or \textit{confounder proxy}. In every role the obstacle is the same --- text is causally relevant only through an unobserved low-dimensional representation, and any analysis that \textit{discovers} that representation from the same sample used for estimation risks representation-induced overfitting. Their remedy is a split-sample workflow: learn the representation on a held-out discovery sample, estimate effects on a disjoint inference sample. \citet{sridhar2022commentary} sharpen the point: causal sufficiency of an embedding is an identification claim, not a modeling choice. \citet{veitch2020adapting} formalize this: an embedding $\hat{Z}$ is causally sufficient if it blocks all backdoor paths through the true latent $Z$; joint fine-tuning to predict treatment and outcome achieves this asymptotically, conditional on no remaining unmeasured confounding. \citet{feder2022causal} survey the full NLP--causality landscape across these three roles.

\textit{The LLM turn.} \citet{imai2024gpi} bypass representation learning entirely: the LLM's \textit{internal} representation is taken as the causal space, and counterfactual texts are generated by steering the model in that space with nuisance features held fixed. Their GPI estimator outperforms observational baselines. \citet{feldman2026sae} decompose LLM activations with sparse autoencoders (SAEs) for an end-to-end pipeline: identify which SAE features encode the treatment construct, steer to generate counterfactuals, then estimate causal effects robustly.

\section*{Critique}

The LLM turn substitutes the ``no unmeasured confounding conditional on $Z$'' assumption with the assumption that the LLM representation is \textit{causally sufficient for the target domain} --- equally untestable in general, and potentially worse if pretraining data under-represents domain-specific institutional context.

The invariance literature converts this into a diagnostic. If causal sufficiency holds, causal effect estimates derived from the LLM representation should be \textit{invariant} across subpopulations defined by stylistically irrelevant variation. Failure of this test --- formalizable via IRM \citep{arjovsky2019irm} or Peters et al.\ \citep{peters2016icp} --- is direct evidence of conflation between causally relevant and irrelevant features. Passing the test is not sufficient for identification, but it is a falsifiable necessary condition these papers do not explicitly derive.

The BO connection is not merely analogical: the SAE feature space of \citet{feldman2026sae} is a low-dimensional continuous action space, and Thompson sampling over it is a principled adaptive design policy for discovering which latent textual features maximally shift the outcome --- a direct instantiation of causal BO \citep{aglietti2021dcbo}.

\section*{Questions}
\begin{enumerate}
    \item Egami et al.'s split-sample workflow halves effective sample size. Is there a Bayesian analog --- a posterior over text representations propagated into the causal estimand --- that maintains calibrated uncertainty without this sacrifice? The marginal likelihood over representation hyperparameters (the GP empirical Bayes analogy from Week 11) seems like the natural starting point.
    \item The invariance test provides a \textit{diagnostic} for LLM causal sufficiency. Does it also yield a training \textit{objective} --- fine-tune the LLM to minimize IRM loss across stylistic environments rather than predictive loss --- and if so, is this strictly stronger than the Veitch et al.\ joint fine-tuning approach or equivalent under their generative model assumptions?
\end{enumerate}

\newpage
\vspace{1em}
\hrule
\vspace{1em}
\section*{Project Update: Text as Confounder Proxy and Treatment in the Zoning Data}
\thispagestyle{empty}

\begin{abstract}
The Austin zoning dataset contains public testimony --- oral and written comment --- associated with each of the 7,074 cases. This week I formalize how text enters the causal model in two distinct roles and map each role to the appropriate method from the readings.
\end{abstract}

\section*{Two Roles for Text in the Zoning Data}

\textit{Role 1: Text as confounder proxy.}
A protest petition signals community opposition, but its \textit{content} reveals the community's mobilization capacity, political framing, and institutional sophistication --- latent confounders for both the petition-filing decision and the council's vote. Standard covariate adjustment cannot condition on raw petition text. The \citet{veitch2020adapting} framework applies: learn a causally sufficient embedding $\hat{Z} = g(\text{petition})$ by jointly fine-tuning a pretrained language model to predict (a) whether the case was appealed post-decision and (b) final vote outcome, then include $\hat{Z}$ as a covariate in the GP surrogate for the objective function $f$.

\textit{Role 2: Text as treatment.}
Whether protest rhetoric appeals to environmental, fiscal, or racial-equity concerns may have a direct causal effect on vote outcome, independent of whether a formal protest petition was filed. The \citet{egami2022texts} split-sample workflow applies: (a) use the pre-2022 discovery sample ($n = 4{,}312$) to learn a topic model or cluster the petition text into rhetorical frames; (b) use the post-2022 inference sample ($n = 2{,}762$) to estimate the causal effect of frame on vote, conditioning on $\hat{Z}$ from Role 1. This separation prevents representation-induced overfitting.

\section*{LLM Augmentation for the Post-2022 Sample}

The change-point identified in previous weeks (2022 municipal election) compounds the text problem: the distribution of petition rhetoric and the council's sensitivity to it may have shifted simultaneously. The post-2022 inference sample is already small ($n = 2{,}762$) and cannot support reliable topic discovery. The \citet{imai2024gpi} GPI approach is attractive here: rather than discovering the rhetorical frame from observational text, use an LLM to \textit{generate} counterfactual petitions with the frame feature varied (environmental vs.\ fiscal framing) and confounders held fixed, then estimate the causal effect from the generated pairs.

The required assumption --- that the LLM's internal representation captures Austin zoning-specific institutional semantics --- is testable via the invariance diagnostic proposed in my critique above: estimate the causal effect of rhetorical frame separately for cases decided by different council committees, and test whether estimates are invariant across committees. Failure of invariance would indicate that the LLM representation conflates framing with committee-specific procedural norms, motivating IRM-style fine-tuning before the GPI estimation step.

\section*{Next Steps}
(1)~Embed petition text using a domain-adapted model (BERT fine-tuned on municipal documents) and perform the invariance test across council committees. (2)~Implement the Egami et al.\ split-sample design for the frame-as-treatment analysis. (3)~If the invariance test fails, implement IRM fine-tuning of the embedding as a preprocessing step. (4)~Integrate $\hat{Z}$ as a covariate in the GP surrogate from previous weeks and evaluate the change in counterfactual comparison results.

\vspace{-0.5em}
{\scriptsize
\setlength{\bibsep}{0pt}
\bibliographystyle{plainnat}
\bibliography{references}
}
\end{document}
